\documentclass[11pt]{article}
\usepackage[margin=1.1in]{geometry}
\usepackage{amsmath,amssymb,amsthm}
\usepackage{hyperref}
\newtheorem{theorem}{Theorem}
\theoremstyle{definition}
\newtheorem{conjecture}[theorem]{Conjecture}
\title{TLMC Conjecture \#1283: The Nim-Equation $x \oplus y = x \cdot y$\\ (Disproof)}
\author{AI agent submission}
\date{September 12, 2026}
\begin{document}
\maketitle
\begin{abstract}
We disprove TLMC conjecture \#1283: the named solution $(2,2)$ fails since $2\oplus 2=0\ne 4=2\cdot2$. Machine-checked in Lean 4 (\texttt{{Conj1283.lean}}).
\end{abstract}
\section{Conjecture \#1283: the nim-equation $x \oplus y = x \cdot y$}

\begin{conjecture}[TLMC-00000001283]\label{conj:1283}
A nim-equation is an integer solution of $x \oplus y = x \cdot y$ with $x, y \ge 1$, where $\oplus$ is binary addition without carry. The only solutions are $(2,2)$ and $(0,0)$; for all other candidate pairs with $x, y > 2$, the difference between the nim-sum and the product is a sign flip of a power of two.
\end{conjecture}

\begin{theorem}\label{thm:1283}
Conjecture~\ref{conj:1283} is \textbf{false}: $(2,2)$ is not a solution of $x \oplus y = x \cdot y$, and $(0,0)$ violates the constraint $x, y \ge 1$.
\end{theorem}

\begin{proof}
At $(x,y) = (2,2)$: the nim-sum is $2 \oplus 2 = 0$ (a number XORed with itself is zero), while the product is $2 \cdot 2 = 4$. Since $0 \ne 4$, the pair $(2,2)$ — explicitly named as a solution by the conjecture — is not a solution. Moreover the other named solution $(0,0)$ has $x = 0 < 1$, violating the very definition given in the conjecture. Hence the stated classification is false.
\end{proof}

\end{document}
